\capitulo{3}{Conceptos teóricos}

\section{Catálogo centralizado y fuentes oficiales}

Una plataforma de descarga puede almacenar los binarios o limitarse a mantener información para encontrarlos. Batch Downloader seguirá el segundo enfoque: conservará metadatos y reglas de resolución, mientras que los instaladores se obtendrán bajo demanda. Esta separación reduce la duplicación de ejecutables y evita que el catálogo se convierta en una fuente desactualizada respecto al editor original.

El concepto de \emph{fuente oficial} no se limita a la URL visible de una página. Una descarga puede atravesar redirecciones o terminar en una red de distribución de contenidos. Por ello, cada aplicación tendrá una lista de dominios permitidos y la URL final se comprobará antes de aceptar el archivo. La interfaz mostrará tanto la página oficial como el estado y la fecha de la última comprobación disponible.

\section{Bundles y personalización temporal}

Un bundle es una colección identificable de aplicaciones. Puede representar una necesidad concreta, por ejemplo ``herramientas para desarrollo web'', o ser simplemente una selección personal. El modelo distingue:

\begin{description}
	\item[Bundle oficial:] creado y mantenido por un administrador.
	\item[Bundle público:] creado por un usuario registrado y visible mediante un enlace compartible.
	\item[Bundle privado:] accesible únicamente para su propietario.
	\item[Bundle temporal:] selección no persistente utilizada para una descarga concreta.
\end{description}

La personalización temporal evita confundir la composición publicada con las preferencias de quien la descarga. Al abrir un bundle público, cualquier usuario podrá añadir o retirar aplicaciones antes de generar su ZIP. Esos cambios vivirán solamente en la solicitud actual. La edición persistente seguirá reservada al propietario o, en el caso de bundles oficiales, a un administrador.

\section{Resolución de enlaces y scraping}

Las URL de descarga cambian por versiones, regiones, sistemas operativos o decisiones del proveedor. Un \emph{resolver} es el componente que transforma la información de una aplicación y el contexto solicitado en una URL descargable validada.

La estrategia prevista será progresiva:

\begin{enumerate}
	\item utilizar una URL directa configurada y todavía válida;
	\item utilizar un resolver específico para el proveedor;
	\item navegar por la página oficial mediante Playwright, inspeccionando enlaces, botones, redirecciones y tráfico de red;
	\item marcar la fuente para revisión manual si no se obtiene un resultado seguro.
\end{enumerate}

Playwright permite controlar navegadores modernos desde Java y ejecutarlos en modo sin interfaz~\cite{playwright}. Su uso será apropiado para páginas que generan los enlaces mediante JavaScript o requieren interacción. La navegación automatizada no elimina las comprobaciones posteriores: el resultado deberá seguir utilizando HTTPS, pertenecer a un dominio permitido y cumplir los límites establecidos.

\section{Procesamiento asíncrono}

Resolver varias páginas, descargar archivos grandes y comprimirlos puede superar el tiempo razonable de una petición HTTP. El procesamiento asíncrono separa la creación de la solicitud de su ejecución.

El backend validará la selección y publicará un mensaje en RabbitMQ. Un worker consumirá ese mensaje y actualizará el estado del trabajo a medida que avance. Spring AMQP proporciona abstracciones para enviar y recibir mensajes y para implementar consumidores dirigidos por mensajes~\cite{springamqp}. El usuario consultará el progreso mediante la API y descargará el ZIP cuando el estado sea \emph{ready}.

Los estados previstos para un trabajo son:

\begin{center}
\begin{tabular}{ll}
\toprule
\textbf{Estado} & \textbf{Significado}\\
\midrule
\emph{queued} & Espera de un worker disponible.\\
\emph{resolving} & Resolución y validación de fuentes.\\
\emph{downloading} & Descarga temporal de instaladores.\\
\emph{compressing} & Construcción del ZIP y del manifiesto.\\
\emph{ready} & Archivo preparado para su entrega.\\
\emph{failed} & No se ha podido producir un resultado válido.\\
\emph{cancelled} & Trabajo cancelado por el usuario o el sistema.\\
\emph{expired} & Resultado eliminado al superar su tiempo de vida.\\
\bottomrule
\end{tabular}
\end{center}

\section{Generación temporal del ZIP}

Cada trabajo utilizará un espacio aislado identificado por su código. Los instaladores no se incorporarán a un almacenamiento permanente del catálogo. MinIO podrá emplearse como almacenamiento de objetos temporal para compartir artefactos entre workers, aplicar tiempos de expiración y evitar que el servidor web mantenga archivos grandes en su sistema local.

El ZIP incluirá una carpeta con los instaladores disponibles, un fichero \ruta{manifest.json} y un documento de lectura. El manifiesto registrará la aplicación, la fuente oficial, el sistema operativo, la arquitectura, el nombre del archivo y el resultado de la validación. Si alguna descarga falla, se añadirá una explicación y el ZIP podrá completarse con los archivos restantes.

\section{Embeddings y búsqueda vectorial}

Un embedding es una representación numérica de un elemento en un espacio vectorial. Textos con significado relacionado tienden a ocupar posiciones próximas según una medida de distancia. Para cada aplicación se generará un vector a partir de su nombre, descripción, categorías, sistemas soportados, arquitectura y otros metadatos relevantes.

PostgreSQL con pgvector permite almacenar estos vectores junto con los datos de referencia y realizar búsquedas de vecinos cercanos mediante distancia euclídea, producto interno o similitud coseno~\cite{pgvector}. Ante una consulta, el sistema generará su embedding y recuperará los candidatos más próximos. Después aplicará reglas deterministas: compatibilidad de sistema y arquitectura, estado activo, restricciones de seguridad y disponibilidad.

La búsqueda semántica complementará, pero no sustituirá, la búsqueda por nombre. Las consultas exactas o parciales seguirán resolviéndose de forma convencional porque son más predecibles y económicas. La búsqueda vectorial se utilizará cuando la consulta exprese una intención, cuando no existan coincidencias suficientes o cuando el usuario active explícitamente el modo semántico.

\section{Modelo de lenguaje y generación aumentada con recuperación}

El modelo de lenguaje recibirá únicamente contexto recuperado desde el catálogo, los estados de los trabajos o registros internos autorizados. Este patrón, conocido como generación aumentada con recuperación, reduce la necesidad de confiar en conocimiento no verificable del modelo.

Durante el desarrollo se prevé utilizar Groq como proveedor de inferencia, cuya API ofrece compatibilidad con clientes del ecosistema OpenAI~\cite{groq}. La integración se encapsulará detrás de una interfaz propia para poder cambiar de proveedor o ejecutar un modelo local sin modificar el resto de la aplicación.

Las respuestas se validarán contra esquemas estructurados antes de mostrarse o transformarse en acciones. Las responsabilidades del modelo quedarán limitadas a:

\begin{itemize}
	\item explicar por qué unas aplicaciones se ajustan a una consulta;
	\item sugerir una selección inicial para un bundle;
	\item resumir errores técnicos en un lenguaje comprensible;
	\item ayudar al administrador a clasificar incidencias o detectar entradas similares.
\end{itemize}

El modelo no podrá aprobar dominios, resolver una URL por sí solo, ejecutar descargas, cambiar reglas de scraping, generar un ZIP ni modificar el catálogo sin una operación autorizada y validada por los servicios correspondientes.

\section{Seguridad en la descarga de recursos externos}

El backend recibirá y seguirá URL externas, lo que introduce riesgo de falsificación de peticiones del lado del servidor o SSRF. Un atacante podría intentar acceder a servicios internos, metadatos de infraestructura o redes privadas. La mitigación prevista incluye:

\begin{itemize}
	\item admitir exclusivamente HTTPS;
	\item resolver DNS y rechazar direcciones locales, privadas, de enlace local o reservadas;
	\item repetir la validación después de cada redirección;
	\item aceptar solo dominios incluidos en la configuración de la aplicación;
	\item limitar redirecciones, tiempo, tamaño y concurrencia;
	\item no ejecutar ningún archivo y sanitizar su nombre antes de incorporarlo al ZIP.
\end{itemize}

También se aplicarán controles de autenticación, autorización, limitación de frecuencia, cuotas y auditoría. Estas medidas evitan que la plataforma se utilice como proxy de descargas masivas y proporcionan evidencias para investigar incidencias.
